\chapter{이 문서에 대하여}

\section{누구를 위한 문서인가}

이 문서는 프로그래밍은 해 보았지만 FPGA나 디지털 회로 설계는 처음인 사람을 위해 썼다. 목표는 Project
Christopher의 코드를 읽고 ``이 코드가 왜 이렇게 생겼고 어떻게 동작하는가''를 이해하는 것이다.
설치와 실행 방법을 안내하는 문서가 아니다. 그 내용은 부록~\ref{app:run}에 짧게만 있다.

Project Christopher는 제2차 세계대전 때 에니그마 암호를 푸는 데 쓰인 기계 \term{봄베}{Bombe}를
재현하는 프로젝트다. 최종 목표는 드럼이 실제로 돌아가는 기계식 실물이고, 전시해서 관람객이 직접 조작해
볼 수 있는 수준을 설계 요구사항으로 둔다. 전시는 같은 봄베를 세 방식으로 만든 세 대로 구성한다.
기어로 드럼을 돌리는 기계식, 릴레이식, 그리고 움직이는 부품이 없는 전자식(FPGA)이다. 나란히 두면 같은
계산이 기계에서 전기기계로, 다시 전자로 옮겨 온 과정이 보인다. 프로그램(Rust)은 세 대 모두가 내야 할
답을 정하는 기준 모델이고, 디지털 회로(SystemVerilog로 기술해 FPGA에 올린다)는 세 번째 전시물이자 실물의
결과를 기준과 대조하는 장치다. 이 문서는 프로그램과 디지털 회로의 코드를 다루고, 마지막 부에서 릴레이식
기계의 전기 회로와 그 회로를 만들고 검증하는 코드를 다룬다.

\section{읽는 순서}

\begin{itemize}
\item \textbf{1부}는 회로 설계의 기본 사고방식이다. 코드를 읽기 전에 알아야 할 개념을 여섯 장에 걸쳐
설명한다. 프로그램을 짜던 사람이 회로를 설계할 때 가장 자주 걸려 넘어지는 곳들이다.
\item \textbf{2부}는 Rust로 쓴 소프트웨어 모델의 해설이다. 봄베가 무엇을 계산하는 기계인지를 여기서
익힌다. 회로는 이 모델과 같은 답을 내야 한다. 끝의 두 장은 일을 나누어 돌리는 탐색(\ref{ch:14-rust-search}장)과
이 코어를 파이썬에서 쓰는 법(\ref{ch:15-python-api}장)이다.
\item \textbf{3부}는 그 모델을 회로로 옮긴 기준 RTL의 해설이다. 1부의 개념이 실제 코드에서 어떻게
나타나는지 본다.
\item \textbf{4부}는 실제 FPGA에 올리는 엔진의 해설이다. 같은 계산을 칩에 들어가는
구조로 다시 짠 코드와, 그 과정에서 내린 결정을 따라간다. 이어서 PC가 보드와 이야기하는 호스트 프로그램과
전시 화면(\ref{ch:31-fpga-host}장), 관람객이 쓰는 메뉴 카드를 만드는 도구(\ref{ch:32-menu-cards}장)를 본다.
\item \textbf{5부}는 하드웨어, 곧 릴레이식 봄베와 두 기계가 함께 쓰는 전기 회로다. 전압·전류·다이오드·트랜지스터와
회로도 읽기(\ref{ch:39-electric-basics}장)로 시작해, 릴레이와 접점으로 논리를
만드는 법(\ref{ch:relay}장), 느리고 들쭉날쭉한 부품으로 시간을 정하는 법(\ref{ch:relay-timing}장), 회로도와
기판을 Python 코드로 생성하는 법(\ref{ch:boardgen}장), 그리고 그 보드를 넷리스트 시뮬레이션으로 기준 모델과
대조하는 법(\ref{ch:boardverify}장), 마지막으로 전원 레일·퓨즈·비상 정지와 인터록·접지(\ref{ch:44-power-safety}장)다.
전기를 몰라도 읽을 수 있게 썼고, 2부의 봄베 원리와 1부의
``기준 모델과 대조한다''(\ref{ch:verify}장)만 알면 된다.
\item \textbf{6부}는 이 프로젝트의 최종 목표인 기어식 봄베다. 드럼과 접점, 톱니 비율과 자리올림, 위치 원판이 어떻게
봄베의 계산을 하는지(\ref{ch:50-gear-mechanics}장), 그 기계를 코드로 그리고 간섭을 검사하는
파라메트릭 CAD(\ref{ch:51-gear-cad}장), 모터를 돌리고 시간 신호를 내는 제어 보드(\ref{ch:52-gear-control}장)를 다룬다.
5부의 시간 규칙(\ref{ch:relay-timing}장)을 먼저 읽으면 좋다.
\end{itemize}

봄베의 원리를 이미 알면 2부를 건너뛰고, 회로 설계를 이미 알면 1부를 건너뛰어도 된다. 5--6부는 3--4부와 따로
읽을 수 있다.

\section{설명하는 방식}

개념 하나를 세 단계로 설명한다.

\Why 그 개념이 그렇게 될 수밖에 없는 원리와 메커니즘을 먼저 적는다.

\General 이 프로젝트와 무관한 교과서적인 회로나 일상의 예를 든다. 카운터, 깜박이는 LED, 덧셈기, FIFO,
신호 처리 필터 같은 것들이다.

\Ours 같은 개념이 이 프로젝트의 코드에서 나타나는 곳을 보인다.

코드는 저장소의 실제 코드에서 발췌했다. 지면을 아끼려고 주석 일부와 오류 처리 등을 생략한 곳이 있으며,
생략한 곳은 \code{...}로 표시했다. 코드 안의 주석은 원문대로 영어다.

\section{봄베를 한 쪽으로}\label{sec:bombe-one-page}

코드를 읽으려면 봄베가 무엇을 하는 기계인지는 알아야 한다. 자세한 설명과 출처는 프로젝트 보고서에 있고,
여기서는 뒤에서 쓸 만큼만 요약한다.

\paragraph{에니그마.} 에니그마는 글자 하나를 누를 때마다 다른 치환표로 바꿔 주는 암호기다. 전류는
\emph{플러그보드, 로터 세 개, 반사판, 로터 세 개(역방향), 플러그보드} 순으로 흐른다. 로터와 반사판을 묶어
\term{스크램블러}{scrambler}라 한다. 플러그보드는 글자 쌍을 맞바꾸는 배선판이고, 서로 이어진 두 글자를
\term{스테커}{Stecker} 짝이라고 한다. 로터는 글자를 칠 때마다 한 칸씩 돌아가므로 스크램블러의 치환은
매번 달라진다. 플러그보드는 하루 동안 그대로다.

\paragraph{크립과 메뉴.} 해독자는 평문의 일부를 추측한다. 이것이 \term{크립}{crib}이다. 크립을 암호문에
맞춰 놓으면 ``$k$번째 글자에서 평문 $p$가 암호문 $c$가 되었다''는 관계가 글자 수만큼 생긴다. 글자를
점으로, 이 관계를 선으로 그린 그래프가 \term{메뉴}{menu}다(그림~\ref{fig:g-menu}).

\begin{figure}[ht]
\centering
\begin{tikzpicture}[x=1cm,y=1cm,
    el/.style={font=\scriptsize,fill=white,inner sep=1pt,text=cgray}]
  \node[ltr,fill=corangebg,draw=corange,line width=1pt] (E) at (0,0) {E};
  \node[ltr] (I) at (2,1.0) {I};
  \node[ltr] (W) at (4,0) {W};
  \node[ltr] (A) at (2,-1.0) {A};
  \node[ltr] (L) at (-2,1.0) {L};
  \node[ltr] (T) at (-4,1.0) {T};
  \node[ltr] (U) at (-2,-1.0) {U};
  \node[ltr] (R) at (4,1.9) {R};
  \node[ltr] (H) at (6,0) {H};
  \draw[line width=1.4pt,cblue] (E) -- node[el] {5} (I) -- node[el] {1} (W) -- node[el] {14} (A) -- node[el] {11} (E);
  \draw (E) -- node[el] {2} (L) -- node[el] {4} (T);
  \draw (E) -- node[el] {16} (U);
  \draw (I) -- node[el] {6} (R);
  \draw (W) -- node[el] {10} (H);
  \node[font=\scriptsize,text=corange,anchor=east] at (-0.42,0) {테스트 글자};
  \node[note,text=cblue] at (2,0) {루프};
\end{tikzpicture}
\caption{메뉴의 예(일부). 선의 숫자는 크립 안에서의 위치다. 파란 선은 출발한 글자로 돌아오는 루프다.}
\label{fig:g-menu}
\end{figure}

\paragraph{봄베가 하는 일.} 봄베는 로터의 가능한 위치 $26^3 = 17{,}576$가지를 하나씩 지나가며, 위치마다
``이 위치에서 메뉴의 모든 관계를 동시에 만족하는 플러그보드가 있을 수 있는가''를 판정한다. 방법은 이렇다.

\begin{itemize}
\item 글자마다 26가닥짜리 케이블이 있다. 케이블 \texttt{E}의 \texttt{v}번 선에 전류가 흐른다는 것은
``\texttt{E}의 스테커 짝이 \texttt{V}다''라는 \emph{가설}이 살아 있다는 뜻이다. 선은 모두
$26\times26 = 676$가닥이다.
\item 메뉴의 선마다 스크램블러 하나가 두 글자의 케이블을 잇는다. 이 연결은 ``\texttt{E}의 짝이 $y$이면
\texttt{I}의 짝은 스크램블러가 $y$를 바꾼 글자다''라는 추론에 해당한다.
\item \term{대각선 보드}{diagonal board}는 ``$X$의 짝이 $y$이면 $Y$의 짝은 $x$다''라는 플러그보드의 대칭성을
배선으로 넣은 것이다. 케이블 $X$의 $y$번 선과 케이블 $Y$의 $x$번 선을 잇는다.
\item 테스트 글자 케이블의 선 하나에 전압을 건다. 전류는 이어진 선으로 퍼져 나가며 그 가설의 논리적
귀결을 모두 통전시킨다.
\item 테스트 글자 케이블(\term{테스트 레지스터}{test register})의 26선이 전부 통전되면 26가지 가설이 모두
모순이므로 그 위치는 버린다. 전부 통전되지는 않으면 모순 없는 가설이 남아 있는 것이므로 기계가
\term{정지}{stop}한다.
\end{itemize}

그림~\ref{fig:g-bombe-front}는 1940년대의 영국 봄베를 문헌의 기술에 따라 그린 것이다. 앞면의 드럼 세 개
한 벌이 위에서 말한 스크램블러 하나이고, 뒷면의 잭에 26선 케이블을 꽂아 메뉴를 배선한다. 뒤에 나오는 코드의
자료 구조와 모듈은 이 그림의 부분들에 하나씩 대응한다.

\begin{figure}[p]
\centering
% 영국 3로터 튜링-웰치먼 봄베의 앞면(위)과 뒷면 잭 배치(아래). 직접 그린 도해이며 사진을 옮긴 것이 아니다.
% 근거: docs/reference/bombe-mechanics.md 2, 5, 11절 (us6812report1944, ellsbury1998bombe, jackson-bh10).
% 단위: 1 = 1 ft. 근거가 있는 것: 뱅크 3개 x 유닛 12개 x 드럼 3개(108개), 지시 드럼 3개가 가운데 뱅크
% 오른쪽, 외형 약 7 ft x 6 ft 6 in, 드럼 지름 약 5 in, 뒷면 잭 열 9개의 용도.
% 어림해 그린 것: 드럼 사이 간격과 캐비닛 안에서의 정확한 위치, 지시 드럼의 배열 방향, 바퀴, 잭의 개수와 간격.
\begin{tikzpicture}[x=1.27cm,y=1.27cm,
    drum/.style={circle,draw=black!70,fill=clight,minimum size=5.3mm,inner sep=0pt,line width=0.4pt},
    lab/.style={font=\footnotesize,align=left,anchor=west},
    lead/.style={draw=cgray,line width=0.4pt}]
  % ---------------- 앞면 ----------------
  \draw[line width=0.9pt,fill=white] (0,0) rectangle (7,6.5);
  \foreach \x in {0.5,6.5} { \draw[fill=cgray!40] (\x,-0.12) circle (0.1); }
  % 드럼: 뱅크 3개 x 3줄 x 12열
  \foreach \b/\by in {0/4.55,1/2.8,2/1.05} {
    \fill[cbluebg!50] (0.33,\by-0.25) rectangle (6.02,\by+1.25);
    \foreach \r in {0,1,2} {
      \foreach \c in {0,...,11} {
        \node[drum] at (0.6+\c*0.47,\by+1.0-\r*0.5) {};
        \fill[black!70] (0.6+\c*0.47,\by+1.0-\r*0.5) circle (0.03);
      }
    }
  }
  % 스크램블러 한 벌 강조 (위 뱅크, 셋째 열)
  \draw[corange,line width=1.1pt,rounded corners=2pt] (1.32,4.32) rectangle (1.76,5.78);
  % 지시 드럼 (가운데 뱅크 오른쪽)
  \foreach \r in {0,1,2} { \node[drum,fill=cyellowbg] (ind\r) at (6.5,3.8-\r*0.5) {}; }
  % 치수
  \draw[<->,cgray] (0,6.8) -- node[above,font=\footnotesize,text=cgray] {폭 약 7 ft (2.1 m)} (7,6.8);
  \draw[<->,cgray] (-0.35,0) -- node[above,rotate=90,font=\footnotesize,text=cgray] {높이 약 6 ft 6 in (2.0 m)} (-0.35,6.5);
  % 라벨
  \draw[lead] (6.02,5.3) -- (7.3,5.9);
  \node[lab] at (7.3,5.9) {\textbf{뱅크}(bank) × 3\\유닛 12개씩, 모두 36개};
  \draw[lead,corange] (1.54,5.78) -- (1.9,6.15) -- (6.5,6.15) -- (7.3,4.95);
  \node[lab,text=corange] at (7.3,4.75) {\textbf{스크램블러 한 벌}(유닛)\\세로로 놓인 드럼 3개\\= 플러그보드 없는 에니그마};
  \draw[lead] (ind0.east) -- (7.3,3.5);
  \node[lab] at (7.3,3.5) {\textbf{지시 드럼}(indicator drums) 3개\\정지한 위치를 읽는다};
  \node[lab] at (7.3,2.1) {한 벌 안의 드럼 (위에서 아래로)\\위: 연속 회전 --- 왼쪽 로터\\가운데 --- 가운데 로터\\아래: 가장 느림 --- 오른쪽 로터};
  \node[lab] at (7.3,0.75) {드럼 108개, 지름 약 5 in\\깊이 약 2 ft, 무게 약 1톤};
  \node[font=\small\bfseries,anchor=west] at (0,7.35) {앞면};
  % ---------------- 뒷면 (잭 배치 도식) ----------------
  \begin{scope}[yshift=-6.6cm]
  \node[font=\small\bfseries,anchor=west] at (0,3.75) {뒷면: 메뉴를 배선하는 잭 (배치 도식)};
  \draw[line width=0.9pt,fill=white] (0,0) rectangle (7,3.4);
  \foreach \g in {0,1,2} {
    \draw[cgray,dashed] (0.25+\g*2.25,0.15) rectangle (2.25+\g*2.25,3.25);
    % 대각선 보드 열
    \foreach \j in {0,...,12} { \fill[cblue] (0.55+\g*2.25,0.4+\j*0.215) rectangle ++(0.22,0.12); }
    % commons 열
    \foreach \j in {0,...,7} { \fill[cgreen] (1.14+\g*2.25,0.5+\j*0.34) rectangle ++(0.22,0.2); }
    % 스크램블러 IN/OUT 열 (유닛당 한 쌍, 12쌍)
    \foreach \j in {0,...,11} {
      \fill[corange] (1.68+\g*2.25,0.4+\j*0.235) rectangle ++(0.13,0.13);
      \fill[corange!55] (1.84+\g*2.25,0.4+\j*0.235) rectangle ++(0.13,0.13);
    }
  }
  % 메뉴 케이블 예 (개념)
  \draw[black!75,line width=0.9pt] (1.75,2.8) .. controls (1.2,2.3) and (1.0,1.9) .. (0.66,1.55);
  \draw[black!75,line width=0.9pt] (1.91,2.8) .. controls (2.3,2.0) and (1.6,1.4) .. (1.25,1.28);
  \draw[black!75,line width=0.9pt] (1.25,1.6) .. controls (0.9,1.2) and (0.9,0.9) .. (0.66,0.68);
  \draw[lead,cblue] (0.66,3.0) -- (0.66,3.5) -- (7.3,3.1);
  \node[lab,text=cblue] at (7.3,3.1) {\textbf{대각선 보드}(diagonal board)\\글자마다 26접점 잭, 뱅크마다 한 벌};
  \draw[lead,cgreen] (5.75,2.0) -- (7.3,2.0);
  \node[lab,text=cgreen] at (7.3,2.0) {\textbf{커먼}(commons)\\한 글자에 선이 여럿 모일 때의 병렬 잭};
  \draw[lead,corange] (6.47,1.0) -- (7.3,0.9);
  \node[lab,text=corange] at (7.3,0.9) {\textbf{스크램블러 IN/OUT 잭}\\유닛마다 한 쌍. 26선 케이블로 잇는다};
  \end{scope}
\end{tikzpicture}

\caption{영국 3로터 봄베의 도해(직접 그린 것이며 사진을 옮긴 것이 아니다). 뱅크와 드럼의 수, 지시 드럼의
위치, 뒷면 잭 열의 용도, 외형 치수는 문헌~\cite{us6812report1944,ellsbury1998bombe,jackson-bh10}을
따랐다. 드럼의 정확한 간격과 위치, 지시 드럼의 배열 방향, 잭의 개수와 간격은 어림해 그렸다.}
\label{fig:g-bombe-front}
\end{figure}

정지의 대부분은 우연히 걸러지지 않은 거짓 정지다. 정지한 위치를 검증하고 키 전체를 구하는 뒷단계를
\term{체킹}{checking}이라 한다.

뒤에서 계속 쓸 사실은 하나다. \emph{봄베의 상태는 676가닥의 선이 각각 통전되었는가이고, 봄베의 계산은
전압을 건 선에서 출발해 이어진 선을 모두 찾는 것이다.}
